% GENERATED FILE. Edit canonical lessons, then run npm run generate:books.
% canonical-source-hash: fnv1a64:72a06857fb53cc66
% canonical-lessons: MR-C376-tolne, MR-C376-parakhne, MR-C376-mulyamapan-karne, MR-C376-gunvatta, MR-C376-darja

\chapter{Weighing Things Up}
\label{ch:mr-weighing-things-up}
\hlchaptermodality{376}

\begin{chapteropening}
\textbf{By the end of this chapter:} \emph{I can weigh up choices and judge quality.}

Complete the last lesson of chapter 376: I can weigh up choices and judge quality.
\end{chapteropening}

\section[\texorpdfstring{\mr{तोलणे} (tolṇe)}{तोलणे (tolṇe)}]{\mr{तोलणे} (tolṇe) — to weigh up}

\label{lesson:MR-C376-tolne}



\begin{quote}
Before the new one: say the Marathi for to comprehend, then the Marathi for to practise.
\end{quote}

\subsection*{You'll want to know: \mr{तोलणे}}
\textbf{\mr{तोलणे}} — \emph{tolṇe} — \textquotedblleft{}to weigh up\textquotedblright{}.

\begin{quote}
\mr{निर्णय} \mr{घेण्यापूर्वी} \mr{त्याने} \mr{दोन्ही} \mr{बाजू} \mr{तोलल्या}.
\end{quote}

\begin{quote}
\emph{nirṇay gheṇyāpūrvī tyāne donhī bājū tolalyā.}
\end{quote}

\begin{quote}
Before deciding, he weighed up both sides.
\end{quote}

Past, with -\mr{ने} (he …): \mr{त्याने} \mr{तोललं}, \emph{tyāne tollaṁ}.

Also to weigh on scales.

\begin{cousinweb}
\textbf{\mr{तोलणे}} (\emph{tolṇe}) is inherited: CDIAL lists Marathi \emph{tol}, \textquotedblleft{}a weight for weighing\textquotedblright{}, under Sanskrit \emph{tōlya}, \textquotedblleft{}to be weighed\textquotedblright{} (CDIAL 5980), and the verb is built on that noun.
\end{cousinweb}

\begin{grammarlens}[title={What you've built}]
One more word for giving reasons and opinions.
\end{grammarlens}

\subsection*{Your turn}
\begin{itemize}
\raggedright
  \item \emph{Say it:} \emph{tolṇe}
  \item \emph{Say it:} \emph{tolṇe}, once more
  \item \emph{Say it:} \emph{nirṇay gheṇyāpūrvī tyāne donhī bājū tolalyā.} — Before deciding, he weighed up both sides.
  \item \emph{Recall:} say the Marathi for to comprehend, then the Marathi for to practise, then say \emph{tolṇe} again
\end{itemize}

\begin{tcolorbox}[breakable,colback=teal!4,colframe=teal!35!black,title={Before you move on}]
What does \mr{तोलणे} mean? (\textquotedblleft{}To weigh up\textquotedblright{}.) Say it once more.
\end{tcolorbox}
\section[\texorpdfstring{\mr{पारखणे} (pārakhṇe)}{पारखणे (pārakhṇe)}]{\mr{पारखणे} (pārakhṇe) — to assess; to judge the worth of}

\label{lesson:MR-C376-parakhne}



\begin{quote}
Before the new one: say the Marathi for to practise, then the Marathi for to weigh up.
\end{quote}

\subsection*{You'll want to know: \mr{पारखणे}}
\textbf{\mr{पारखणे}} — \emph{pārakhṇe} — \textquotedblleft{}to assess; to judge the worth of\textquotedblright{}.

\begin{quote}
\mr{माणसं} \mr{पारखायला} \mr{वेळ} \mr{लागतो}.
\end{quote}

\begin{quote}
\emph{māṇsaṁ pārkhāylā veḷ lāgto.}
\end{quote}

\begin{quote}
It takes time to judge people.
\end{quote}

Past, with -\mr{ने} (he …): \mr{त्याने} \mr{पारखलं}, \emph{tyāne pārakhlaṁ}.

\mr{पारख}, the eye of a judge: \mr{त्याला} \mr{सोन्याची} \mr{पारख} \mr{आहे}.

\begin{cousinweb}
\textbf{\mr{पारखणे}} (\emph{pārakhṇe}) is a tadbhava: CDIAL lists Marathi \emph{pārakhṇẽ}, \textquotedblleft{}to examine\textquotedblright{}, under Sanskrit \emph{parīkṣate}, \textquotedblleft{}observes, examines\textquotedblright{} (CDIAL 7904); \mr{परीक्षा}, an exam, is its learned cousin.
\end{cousinweb}

\begin{grammarlens}[title={What you've built}]
One more word for giving reasons and opinions.
\end{grammarlens}

\subsection*{Your turn}
\begin{itemize}
\raggedright
  \item \emph{Say it:} \emph{pārakhṇe}
  \item \emph{Say it:} \emph{pārakhṇe}, once more
  \item \emph{Say it:} \emph{māṇsaṁ pārkhāylā veḷ lāgto.} — It takes time to judge people.
  \item \emph{Recall:} say the Marathi for to practise, then the Marathi for to weigh up, then say \emph{pārakhṇe} again
\end{itemize}

\begin{tcolorbox}[breakable,colback=teal!4,colframe=teal!35!black,title={Before you move on}]
What does \mr{पारखणे} mean? (\textquotedblleft{}To assess; to judge the worth of\textquotedblright{}.) Say it once more.
\end{tcolorbox}
\section[\texorpdfstring{\mr{मूल्यमापन} \mr{करणे} (mūlyamāpan karṇe)}{मूल्यमापन करणे (mūlyamāpan karṇe)}]{\mr{मूल्यमापन} \mr{करणे} (mūlyamāpan karṇe) — to evaluate}

\label{lesson:MR-C376-mulyamapan-karne}



\begin{quote}
Before the new one: say the Marathi for to weigh up, then the Marathi for to assess.
\end{quote}

\subsection*{You'll want to know: \mr{मूल्यमापन} \mr{करणे}}
\textbf{\mr{मूल्यमापन} \mr{करणे}} — \emph{mūlyamāpan karṇe} — \textquotedblleft{}to evaluate\textquotedblright{}.

\begin{quote}
\mr{शिक्षक} \mr{वर्षभर} \mr{मुलांच्या} \mr{कामाचं} \mr{मूल्यमापन} \mr{करतात}.
\end{quote}

\begin{quote}
\emph{śikṣak varṣabhar mulāñcyā kāmācaṁ mūlyamāpan kartāt.}
\end{quote}

\begin{quote}
Teachers evaluate the children's work all year round.
\end{quote}

Past, with -\mr{ने} (he …): \mr{त्याने} \mr{मूल्यमापन} \mr{केलं}, \emph{tyāne mūlyamāpan kelaṁ}.

What is evaluated takes -\mr{चं}: \mr{कामाचं} \mr{मूल्यमापन}.

\begin{cousinweb}
\textbf{\mr{मूल्यमापन} \mr{करणे}} (\emph{mūlyamāpan karṇe}) joins \mr{मूल्यमापन}, a learned compound built from Sanskrit \emph{mūlya}, \textquotedblleft{}value\textquotedblright{}, and \emph{māpana}, \textquotedblleft{}measuring\textquotedblright{} (Monier-Williams), to \mr{करणे}, \textquotedblleft{}to do\textquotedblright{}.
\end{cousinweb}

\begin{grammarlens}[title={What you've built}]
One more word for giving reasons and opinions.
\end{grammarlens}

\subsection*{Your turn}
\begin{itemize}
\raggedright
  \item \emph{Say it:} \emph{mūlyamāpan karṇe}
  \item \emph{Say it:} \emph{mūlyamāpan karṇe}, once more
  \item \emph{Say it:} \emph{śikṣak varṣabhar mulāñcyā kāmācaṁ mūlyamāpan kartāt.} — Teachers evaluate the children's work all year round.
  \item \emph{Recall:} say the Marathi for to weigh up, then the Marathi for to assess, then say \emph{mūlyamāpan karṇe} again
\end{itemize}

\begin{tcolorbox}[breakable,colback=teal!4,colframe=teal!35!black,title={Before you move on}]
What does \mr{मूल्यमापन} \mr{करणे} mean? (\textquotedblleft{}To evaluate\textquotedblright{}.) Say it once more.
\end{tcolorbox}
\section[\texorpdfstring{\mr{गुणवत्ता} (guṇvattā)}{गुणवत्ता (guṇvattā)}]{\mr{गुणवत्ता} (guṇvattā) — quality}

\label{lesson:MR-C376-gunvatta}



\begin{quote}
Before the new one: say the Marathi for to assess, then the Marathi for to evaluate.
\end{quote}

\subsection*{You'll want to know: \mr{गुणवत्ता}}
\textbf{\mr{गुणवत्ता}} — \emph{guṇvattā} — \textquotedblleft{}quality\textquotedblright{}.

\begin{quote}
\mr{या} \mr{दुकानातल्या} \mr{भाज्यांची} \mr{गुणवत्ता} \mr{चांगली} \mr{आहे}.
\end{quote}

\begin{quote}
\emph{yā dukānātalyā bhājyāñcī guṇvattā cāṅglī āhe.}
\end{quote}

\begin{quote}
The quality of the vegetables in this shop is good.
\end{quote}

Feminine.

\mr{गुण}, merit, + -\mr{वत्ता}.

\begin{cousinweb}
\textbf{\mr{गुणवत्ता}} (\emph{guṇvattā}) is a tatsama, taken from Sanskrit \emph{guṇavattā}, \textquotedblleft{}excellence, possessing merit\textquotedblright{} (Monier-Williams).
\end{cousinweb}

\begin{grammarlens}[title={What you've built}]
One more word for giving reasons and opinions.
\end{grammarlens}

\subsection*{Your turn}
\begin{itemize}
\raggedright
  \item \emph{Say it:} \emph{guṇvattā}
  \item \emph{Say it:} \emph{guṇvattā}, once more
  \item \emph{Say it:} \emph{yā dukānātalyā bhājyāñcī guṇvattā cāṅglī āhe.} — The quality of the vegetables in this shop is good.
  \item \emph{Recall:} say the Marathi for to assess, then the Marathi for to evaluate, then say \emph{guṇvattā} again
\end{itemize}

\begin{tcolorbox}[breakable,colback=teal!4,colframe=teal!35!black,title={Before you move on}]
What does \mr{गुणवत्ता} mean? (\textquotedblleft{}Quality\textquotedblright{}.) Say it once more.
\end{tcolorbox}
\section[\texorpdfstring{\mr{दर्जा} (darjā)}{दर्जा (darjā)}]{\mr{दर्जा} (darjā) — a standard; a grade}

\label{lesson:MR-C376-darja}



\begin{quote}
Before the new one: say the Marathi for to evaluate, then the Marathi for quality.
\end{quote}

\subsection*{You'll want to know: \mr{दर्जा}}
\textbf{\mr{दर्जा}} — \emph{darjā} — \textquotedblleft{}a standard; a grade\textquotedblright{}.

\begin{quote}
\mr{या} \mr{शाळेचा} \mr{दर्जा} \mr{खूप} \mr{वरचा} \mr{आहे}.
\end{quote}

\begin{quote}
\emph{yā śāḷecā darjā khūp varcā āhe.}
\end{quote}

\begin{quote}
This school's standard is very high.
\end{quote}

Masculine.

\mr{दर्जेदार}, of good quality.

\begin{cousinweb}
\textbf{\mr{दर्जा}} (\emph{darjā}) is a loan from Persian, usually traced to Arabic \emph{daraja}, \textquotedblleft{}step, rank\textquotedblright{}.
\end{cousinweb}

\begin{grammarlens}[title={What you've built}]
One more word for giving reasons and opinions.
\end{grammarlens}

\subsection*{Your turn}
\begin{itemize}
\raggedright
  \item \emph{Say it:} \emph{darjā}
  \item \emph{Say it:} \emph{darjā}, once more
  \item \emph{Say it:} \emph{yā śāḷecā darjā khūp varcā āhe.} — This school's standard is very high.
  \item \emph{Recall:} say the Marathi for to evaluate, then the Marathi for quality, then say \emph{darjā} again
\end{itemize}

\begin{tcolorbox}[breakable,colback=teal!4,colframe=teal!35!black,title={Before you move on}]
What does \mr{दर्जा} mean? (\textquotedblleft{}A standard; a grade\textquotedblright{}.) Say it once more.
\end{tcolorbox}
